\begin{derivation}[从\eqnrefs{eq:generic-vector-ff-exact-width-dp92}到\eqnrefs{eq:zff-width-tree-dp11}]
令
\begin{equation*}
 q_V\equiv\frac{M_V}{c},\qquad a\equiv m_1c,\qquad b\equiv m_2c,
\end{equation*}
于是 $P^2=q_V^2$、$p^2=a^2$、$k^2=b^2$。先把手征顶角改写为矢量--轴矢量形式
\begin{equation*}
 \Gamma^\mu
 =\gamma^\mu(g_V-g_A\gamma^5),\qquad
 g_V=\frac{c_L+c_R}{2},\qquad
 g_A=\frac{c_L-c_R}{2}.
\end{equation*}
对末态自旋求和得到
\begin{align*}
 T^{\mu\nu}
 &=\Tr\!\left[(\slashed p+a)\gamma^\mu(g_V-g_A\gamma^5)
 (\slashed k-b)\gamma^\nu(g_V^*-g_A^*\gamma^5)\right].
\end{align*}
极化投影张量关于 $\mu,\nu$ 对称，因此含 $\epsilon^{\mu\nu\alpha\beta}$ 的矢量--轴矢量交叉项不贡献总宽度。两个对称迹分别为
\begin{align*}
 T_{VV}^{\mu\nu}
 &=4\left[p^\mu k^\nu+p^\nu k^\mu
 -\eta^{\mu\nu}(p\!\cdot\!k+ab)\right],\\
 T_{AA}^{\mu\nu}
 &=4\left[p^\mu k^\nu+p^\nu k^\mu
 -\eta^{\mu\nu}(p\!\cdot\!k-ab)\right].
\end{align*}
利用
\begin{equation*}
 |g_V|^2+|g_A|^2=\frac{|c_L|^2+|c_R|^2}{2},\qquad
 |g_A|^2-|g_V|^2=-\operatorname{Re}(c_Lc_R^*),
\end{equation*}
可把对称部分整理为
\begin{equation*}
 T_{\rm sym}^{\mu\nu}
 =2S\left(p^\mu k^\nu+p^\nu k^\mu-\eta^{\mu\nu}p\!\cdot\!k\right)
 -4R\,ab\,\eta^{\mu\nu},
\end{equation*}
其中
\begin{equation*}
 S\equiv|c_L|^2+|c_R|^2,\qquad
 R\equiv\operatorname{Re}(c_Lc_R^*).
\end{equation*}
有质量矢量场的三个物理极化给出
\begin{equation*}
 \Pi_{\mu\nu}(P)
 =-\eta_{\mu\nu}+\frac{P_\mu P_\nu}{q_V^2}.
\end{equation*}
由二体在壳关系
\begin{equation*}
 p\!\cdot\!k=\frac{q_V^2-a^2-b^2}{2},\qquad
 P\!\cdot\!p=\frac{q_V^2+a^2-b^2}{2},\qquad
 P\!\cdot\!k=\frac{q_V^2-a^2+b^2}{2},
\end{equation*}
先算动量张量的投影：
\begin{align*}
 &\Pi_{\mu\nu}
 \left(p^\mu k^\nu+p^\nu k^\mu-\eta^{\mu\nu}p\!\cdot\!k\right)\\
 &\quad=p\!\cdot\!k
 +\frac{2(P\!\cdot\!p)(P\!\cdot\!k)}{q_V^2}\\
 &\quad=q_V^2-\frac{a^2+b^2}{2}
 -\frac{(a^2-b^2)^2}{2q_V^2}.
\end{align*}
同时
\begin{equation*}
 \Pi_{\mu\nu}\eta^{\mu\nu}=-4+1=-3.
\end{equation*}
因此对初态三个极化平均后的振幅平方为
\begin{align*}
 \overline{|\mathcal M|^2}
 &=\frac13\Pi_{\mu\nu}T_{\rm sym}^{\mu\nu}\\
 &=\frac{q_V^2}{3}
 \left\{
 S\left[2-r_1-r_2-(r_1-r_2)^2\right]
 +12R\sqrt{r_1r_2}
 \right\},
\end{align*}
其中 $r_i=m_i^2c^4/M_V^2=a_i^2/q_V^2$。

精确二体 Lorentz 不变相空间在母粒子静止系为
\begin{align*}
 \int\dd\Phi_2
 &=\frac{1}{8\pi}
 \frac{\sqrt{\lambda(q_V^2,a^2,b^2)}}{q_V^2}\\
 &=\frac{\sqrt{\lambda(1,r_1,r_2)}}{8\pi}.
\end{align*}
将它和上式代回正文\eqnrefs{eq:generic-vector-ff-exact-width-dp92}，并用
$q_V^2=M_V^2/c^2$，得到
\begin{equation*}
 \Gamma_{V\to12}
 =N_c^f\frac{M_V}{48\pi\hbar}\sqrt{\lambda(1,r_1,r_2)}
 \left\{
 S\left[2-r_1-r_2-(r_1-r_2)^2\right]
 +12R\sqrt{r_1r_2}
 \right\},
\end{equation*}
即正文\eqnrefs{eq:generic-vector-ff-exact-closed-dp92}。

正文常用的轻费米子公式定义为有限质量闭式在原点的连续极限。由于
\begin{equation*}
 \lim_{(r_1,r_2)\to(0,0)}\sqrt{\lambda(1,r_1,r_2)}=1,
 \qquad
 \lim_{(r_1,r_2)\to(0,0)}\sqrt{r_1r_2}=0,
\end{equation*}
把 $r_1=r_2=0$ 代入前面的有限质量闭式，其方括号化为 $2S$，从而
\begin{align*}
\lim_{(m_1,m_2)\to(0,0)}\Gamma_{V\to12}
&=N_c^f\frac{M_V}{48\pi\hbar}(2S)\\
&=N_c^f\frac{M_V}{24\pi\hbar}S.
\end{align*}
即
\begin{equation*}
 \boxed{
 \Gamma_{V\to f\bar f}^{(m_f=0)}
 \equiv
 \lim_{(m_1,m_2)\to(0,0)}\Gamma_{V\to12}
 =N_c^f\frac{M_V}{24\pi\hbar}S.}
\end{equation*}
这就是正文\eqnrefs{eq:generic-vector-ff-width-dp91} 所使用的无质量末态子模型；任意有限末态质量均由前面的精确闭式而不是由该极限式描述。

在无质量末态子模型中，对 $W\to\ell\nu$ 代入 $N_c^f=1$、$c_L=g/\sqrt2$、$c_R=0$，得到
\begin{equation*}
 \Gamma(W\to\ell\nu)=\frac{g^2M_W}{48\pi\hbar}.
\end{equation*}
再用 $g=2M_W/v_E$、$v_E^{-2}=\sqrt2G_F^{(E)}$，得到正文\eqnrefs{eq:w-leptonic-width-dp11}。有限 $m_\ell,m_\nu$ 时应回到\eqnrefs{eq:generic-vector-ff-exact-closed-dp92}。同样，在无质量末态子模型中对 $Z\to f\bar f$，
\begin{equation*}
 |c_L^Z|^2+|c_R^Z|^2
 =\frac{g^2}{2c_W^2}\left[(g_V^f)^2+(g_A^f)^2\right],
\end{equation*}
故
\begin{equation*}
 \Gamma(Z\to f\bar f)
 =N_c^f\frac{g^2M_Z}{48\pi c_W^2\hbar}
 \left[(g_V^f)^2+(g_A^f)^2\right].
\end{equation*}
最后用 $M_W=c_WM_Z$ 与\eqnrefs{eq:gf-matching-dp11} 得到正文\eqnrefs{eq:zff-width-tree-dp11}。
\end{derivation}
